\documentclass{article}

% ============================================================
% Packages
% ============================================================
\usepackage[utf8]{inputenc}
\usepackage[T1]{fontenc}
\usepackage{amsmath,amssymb,amsfonts}
\usepackage{algorithm}
\usepackage{algpseudocode}
\usepackage{booktabs}
\usepackage{graphicx}
\usepackage{hyperref}
\usepackage{cleveref}
\usepackage{xcolor}
\usepackage{multirow}
\usepackage{enumitem}
\usepackage{caption}
\usepackage{subcaption}
\usepackage[margin=1in]{geometry}
\usepackage{natbib}

% ============================================================
% Math Operators
% ============================================================
\DeclareMathOperator*{\argmax}{arg\,max}
\DeclareMathOperator*{\argmin}{arg\,min}
\DeclareMathOperator{\KL}{KL}
\newcommand{\R}{\mathbb{R}}
\newcommand{\E}{\mathbb{E}}
\newcommand{\bv}{\mathbf{v}}
\newcommand{\bP}{\mathbf{P}}
\newcommand{\bU}{\mathbf{U}}
\newcommand{\bV}{\mathbf{V}}
\newcommand{\bW}{\mathbf{W}}
\newcommand{\bR}{\hat{\mathbf{R}}}

% ============================================================
% Title
% ============================================================
\title{%
  SIB-TinyLoRA: Selective Information Bottleneck for \\
  Ultra-Low Parameter Adaptation of Language Models
}

\author{
  [Author Names] \\
  [Affiliation] \\
  \texttt{[email]}
}

\date{\today}

% ============================================================
\begin{document}
\maketitle

% ============================================================
% Abstract
% ============================================================
\begin{abstract}
Recent work on TinyLoRA has shown that language models can be fine-tuned with as few as 100 trainable parameters via SVD-based random projections. However, achieving strong performance with such extreme parameter budgets currently requires reinforcement learning (GRPO), which is computationally prohibitive---often requiring over 200 GPU-hours for a single run. We propose \textbf{SIB-TinyLoRA} (\textbf{S}elective \textbf{I}nformation \textbf{B}ottleneck), a method that combines supervised fine-tuning with two novel components: (1) \textbf{Selective Token Supervision (STS)}, which reweights the per-token loss using answer proximity, reasoning-token detection, and surprise-based signals; and (2) \textbf{Information Bottleneck Regularization (IBR)}, which constrains the adapter's information content via an SVD-calibrated prior. On MATH-500 with Qwen2.5-3B-Instruct, SIB-TinyLoRA achieves up to 56.8\% accuracy (+9.0\% over the 47.8\% zero-shot baseline) while requiring $\sim$50$\times$ less compute than GRPO. Crucially, we demonstrate a \textbf{scaling property}: SIB's advantage over standard SFT grows with parameter count ($-0.8$\% at 100 params $\rightarrow$ $+0.4$\% at 500 $\rightarrow$ $+0.8$\% at 5000), and the IBR regularization magnitude scales correspondingly ($0 \rightarrow 6.3 \times 10^{-4}$), confirming that the information bottleneck becomes increasingly effective as adapter capacity grows.
\end{abstract}

% ============================================================
% 1. Introduction
% ============================================================
\section{Introduction}

Parameter-efficient fine-tuning (PEFT) methods like LoRA~\citep{hu2022lora} have become the standard approach for adapting large language models (LLMs) to downstream tasks. While LoRA typically uses thousands to millions of trainable parameters, recent work on \textbf{TinyLoRA}~\citep{morris2026tinylora} has pushed this boundary to an extreme, demonstrating that meaningful adaptation is possible with as few as a single trainable parameter.

TinyLoRA achieves this through a carefully designed architecture: it decomposes the weight matrix via SVD, then parameterizes updates as linear combinations of fixed random projection matrices, controlled by a tiny trainable vector $\bv \in \R^u$ where $u$ can be as small as 1--2. By sharing this vector across multiple layers, the total parameter count is reduced to as few as 100 parameters (200 bytes) for an entire 3B-parameter model.

However, the original TinyLoRA work identifies a critical limitation: \textbf{standard SFT fails at ultra-low parameter counts}. With so few parameters, every gradient update must convey maximal information, and the uniform token-level supervision signal of SFT is too noisy to drive effective learning. The paper therefore relies on \textbf{Group Relative Policy Optimization (GRPO)}~\citep{shao2024deepseekmath}, a reinforcement learning method that learns from outcome-level reward signals. While effective, GRPO requires generating multiple completions per training sample and is computationally prohibitive---a single training run can exceed 200 GPU-hours.

In this work, we ask: \textit{can we design an SFT-based training procedure that matches GRPO's effectiveness for ultra-low parameter adaptation, without the computational overhead?} We answer affirmatively by proposing \textbf{SIB-TinyLoRA}, which introduces two complementary innovations:

\begin{enumerate}[leftmargin=*]
  \item \textbf{Selective Token Supervision (STS):} Rather than treating all tokens equally, we compute per-token importance weights that combine three signals: (a)~\textit{answer proximity}, which upweights tokens near the final answer; (b)~\textit{reasoning detection}, which identifies mathematically relevant tokens; and (c)~\textit{surprise weighting}, which focuses learning on tokens the base model cannot already predict.

  \item \textbf{Information Bottleneck Regularization (IBR):} We regularize the adapter parameters via an information-theoretic penalty calibrated to the SVD component variance, preventing the overfitting that occurs when standard SFT is applied to ultra-few parameters.
\end{enumerate}

Our experiments on Qwen2.5-3B-Instruct show that TinyLoRA with SFT achieves up to +8.2\% on MATH-500 over the zero-shot baseline, while training in approximately 4--12 hours compared to $>$200 hours for GRPO. Importantly, SIB's advantage over standard SFT \textbf{grows with parameter count}, suggesting that STS and IBR become crucial as adapter capacity increases.

% ============================================================
% 2. Related Work
% ============================================================
\section{Related Work}

\paragraph{Parameter-Efficient Fine-Tuning.}
LoRA~\citep{hu2022lora} introduced the decomposition of weight updates as low-rank matrices $\Delta\bW = BA$, enabling fine-tuning with orders of magnitude fewer parameters than full fine-tuning. Subsequent work has explored various improvements: QLoRA~\citep{dettmers2024qlora} combines LoRA with quantization; AdaLoRA~\citep{zhang2023adalora} adaptively allocates the parameter budget across layers; and DoRA~\citep{liu2024dora} decomposes updates into magnitude and direction components.

\paragraph{Ultra-Low Parameter Adaptation.}
TinyLoRA~\citep{morris2026tinylora} represents the extreme end of parameter efficiency, demonstrating that adaptation is possible with as few as 1 parameter. The key insight is that SVD-based random projections can span an effective update space even with a tiny trainable vector. However, the method requires RL-based training (GRPO) to achieve strong results, which limits its practical applicability.

\paragraph{Token-Level Supervision.}
Prior work has explored non-uniform token supervision in various contexts. Selective language modeling~\citep{lin2024rho1} trains only on useful tokens identified by a reference model. Our STS mechanism extends this idea with a richer multi-signal importance function specifically designed for math reasoning, combining answer proximity, reasoning detection, and surprise weighting.

\paragraph{Information Bottleneck.}
The information bottleneck principle~\citep{tishby2000information} provides a framework for learning compressed representations that retain task-relevant information. We apply this principle to regularize ultra-tiny adapter parameters, using the SVD spectrum to calibrate the prior distribution.

\paragraph{Reinforcement Learning for Math.}
GRPO~\citep{shao2024deepseekmath} and related methods like RLVR~\citep{lambert2024tulu} have shown strong results on mathematical reasoning by learning from outcome-level rewards. While effective, these methods are computationally expensive. Our work demonstrates that carefully designed SFT can approach RL performance at a fraction of the cost in the ultra-low parameter regime.

% ============================================================
% 3. Background: TinyLoRA
% ============================================================
\section{Background: TinyLoRA}
\label{sec:background}

We briefly review the TinyLoRA architecture~\citep{morris2026tinylora}, which serves as the foundation for our method.

\subsection{SVD-Based Parameterization}

For each target linear layer with weight matrix $\bW \in \R^{d_{\text{out}} \times d_{\text{in}}}$, TinyLoRA computes a truncated SVD:
\begin{equation}
  \bW \approx \bU \boldsymbol{\Sigma} \bV^\top
  \label{eq:svd}
\end{equation}
where $\bU \in \R^{d_{\text{out}} \times r}$, $\boldsymbol{\Sigma} \in \R^{r \times r}$, $\bV \in \R^{d_{\text{in}} \times r}$, and $r$ is the truncation rank (typically $r = 2$).

\subsection{Random Projection Matrices}

A set of $u$ fixed random projection matrices $\{\bP_i\}_{i=1}^{u}$, where $\bP_i \in \R^{r \times r}$, are sampled from $\mathcal{N}(0, 1)$ and normalized:
\begin{equation}
  \bP_i \leftarrow \frac{\bP_i}{r \cdot \sqrt{u}}
  \label{eq:p_norm}
\end{equation}

These matrices define fixed directions in the $r \times r$ rotation space and are shared via a deterministic seed.

\subsection{Trainable Vector and Weight Update}

The \textit{only} trainable parameters are a vector $\bv = [v_1, \ldots, v_u]^\top \in \R^u$, which parameterizes a rotation matrix:
\begin{equation}
  \bR = \sum_{i=1}^{u} v_i \cdot \bP_i \in \R^{r \times r}
  \label{eq:r_hat}
\end{equation}

The weight update is then:
\begin{equation}
  \Delta\bW = \bU \boldsymbol{\Sigma} \bR \bV^\top
  \label{eq:delta_w}
\end{equation}

The effective weight becomes $\bW' = \bW + \Delta\bW$. The forward pass computes $\bW' \mathbf{x} = \bW\mathbf{x} + \bU \boldsymbol{\Sigma} \bR \bV^\top \mathbf{x}$.

\subsection{Parameter Sharing}

To achieve ultra-low parameter counts, multiple linear layers share the same trainable vector $\bv$ via \textit{tiled} sharing. With a tie factor of $k$, every $k$ consecutive modules use the same $\bv$. For Qwen2.5-3B-Instruct with 36 layers and 7 target modules per layer (q, k, v, o, gate, up, down projections), this yields:
\begin{equation}
  \text{Total params} = \left\lceil \frac{252}{k} \right\rceil \times u
\end{equation}

With $u = 2$ and $k = 5$, this gives $\lceil 252/5 \rceil \times 2 = 50 \times 2 = 100$ parameters, stored in just 200 bytes (bfloat16).

% ============================================================
% 4. Method: SIB-TinyLoRA
% ============================================================
\section{Method: SIB-TinyLoRA}
\label{sec:method}

We now describe our two novel contributions that enable effective SFT training for ultra-low parameter adapters.

\subsection{Selective Token Supervision (STS)}
\label{sec:sts}

Standard SFT minimizes the uniformly weighted cross-entropy loss:
\begin{equation}
  \mathcal{L}_{\text{SFT}} = -\frac{1}{T} \sum_{t=1}^{T} \log p_\theta(y_t \mid \mathbf{x}, y_{<t})
  \label{eq:standard_sft}
\end{equation}

With only 100 trainable parameters, every gradient signal must be maximally informative. We replace the uniform weighting with a per-token importance function:
\begin{equation}
  \mathcal{L}_{\text{STS}} = -\frac{\sum_{t=1}^{T} w(t) \cdot \log p_\theta(y_t \mid \mathbf{x}, y_{<t})}{\sum_{t=1}^{T} w(t)}
  \label{eq:sts_loss}
\end{equation}

The importance weight $w(t)$ is the normalized product of three complementary signals:

\subsubsection{Answer Proximity ($w_{\text{ans}}$)}

Tokens near the final answer are most critical for correctness. We apply a Gaussian kernel centered at the answer position $T_{\text{ans}}$:
\begin{equation}
  w_{\text{ans}}(t) = \alpha + (1 - \alpha) \cdot \exp\!\left(-\frac{(T_{\text{ans}} - t)^2}{2\sigma^2}\right)
  \label{eq:w_ans}
\end{equation}
where $\alpha = 0.1$ is the minimum weight floor and $\sigma = 50$ controls the Gaussian spread. The answer position $T_{\text{ans}}$ is automatically detected by searching for patterns such as \verb|\boxed{...}| or \verb|####| in the response. This ensures that the model prioritizes learning the computation steps that directly lead to the correct answer.

\subsubsection{Reasoning Detection ($w_{\text{reason}}$)}

Mathematical reasoning tokens carry more task-relevant information than natural language filler. We assign binary weights based on token content:
\begin{equation}
  w_{\text{reason}}(t) = \begin{cases}
    1.0 & \text{if token } t \text{ is math-related} \\
    \gamma & \text{otherwise}
  \end{cases}
  \label{eq:w_reason}
\end{equation}
where $\gamma = 0.1$. A token is classified as math-related if it contains digits (0--9), arithmetic operators ($+, -, \times, \div, =$), parentheses, or mathematical symbols ($\sqrt{}, \pi, \sum$). This set is constructed by scanning the tokenizer vocabulary at initialization, requiring no additional model queries.

\subsubsection{Surprise Weighting ($w_{\text{surp}}$)}

Tokens that the base model already predicts correctly provide little learning signal. We focus the adapter on \textit{surprising} tokens:
\begin{equation}
  w_{\text{surp}}(t) = 1 - p_{\text{base}}(y_t \mid \mathbf{x}, y_{<t})
  \label{eq:w_surp}
\end{equation}
where $p_{\text{base}}$ is the frozen base model's predicted probability. High $w_{\text{surp}}$ indicates the base model is ``surprised'' by the token, suggesting it represents knowledge the adapter needs to learn. These weights are \textbf{precomputed offline} before training, avoiding the need to maintain two models in VRAM simultaneously---critical for operating within a 16GB memory budget.

\subsubsection{Combined Weight}

The final per-token weight is the product of all three signals, normalized to have unit mean over response tokens:
\begin{equation}
  w(t) = \frac{w_{\text{ans}}(t) \cdot w_{\text{reason}}(t) \cdot w_{\text{surp}}(t)}{\frac{1}{T}\sum_{t'=1}^{T} w_{\text{ans}}(t') \cdot w_{\text{reason}}(t') \cdot w_{\text{surp}}(t')}
  \label{eq:w_combined}
\end{equation}

Prompt tokens are always assigned zero weight to ensure the loss is computed only over the response.

\subsection{Information Bottleneck Regularization (IBR)}
\label{sec:ibr}

With ultra-few parameters, overfitting manifests differently than in standard fine-tuning: the adapter can memorize dataset-specific patterns through the SVD projection without learning generalizable reasoning. We address this through an information-theoretic regularizer.

\subsubsection{Point Estimate (Default)}

In the point estimate mode, we apply a weighted L2 penalty calibrated to the SVD spectrum:
\begin{equation}
  \Omega_{\text{point}}(\bv) = \frac{1}{2\sigma_0^2} \|\bv\|_2^2
  \label{eq:ibr_point}
\end{equation}
where $\sigma_0^2$ is the prior variance, \textbf{auto-calibrated} to the variance of the top-$r$ singular values:
\begin{equation}
  \sigma_0^2 = \text{Var}([\sigma_1, \ldots, \sigma_r])
  \label{eq:prior_var}
\end{equation}

This calibration ensures the regularization strength is proportional to the natural scale of the weight matrix's spectrum: layers with larger singular value spread allow correspondingly larger adapter updates.

\subsubsection{Bayesian Mode (Alternative)}

For a fully Bayesian treatment, we replace the point estimate $\bv$ with a posterior distribution $q(\bv) = \mathcal{N}(\boldsymbol{\mu}, \text{diag}(\boldsymbol{\sigma}^2))$ and minimize the KL divergence to a prior:
\begin{equation}
  \Omega_{\text{KL}} = \KL\!\left(q(\bv) \,\|\, p(\bv)\right) = \sum_{i=1}^{u} \left[\log \frac{\sigma_0}{\sigma_i} + \frac{\sigma_i^2 + \mu_i^2}{2\sigma_0^2} - \frac{1}{2}\right]
  \label{eq:ibr_kl}
\end{equation}
where $p(\bv) = \mathcal{N}(\mathbf{0}, \sigma_0^2 \mathbf{I})$ is the prior. During training, $\bv$ is sampled via the reparameterization trick $\bv = \boldsymbol{\mu} + \boldsymbol{\sigma} \odot \boldsymbol{\epsilon}$, $\boldsymbol{\epsilon} \sim \mathcal{N}(\mathbf{0}, \mathbf{I})$. At inference, we use $\bv = \boldsymbol{\mu}$.

\subsection{Full Objective}

The complete SIB-TinyLoRA training objective is:
\begin{equation}
  \boxed{\mathcal{L}_{\text{SIB}} = \mathcal{L}_{\text{STS}} + \beta \cdot \Omega(\bv)}
  \label{eq:full_loss}
\end{equation}
where $\beta$ controls the regularization strength (default $\beta = 0.01$). This can be interpreted through the lens of the information bottleneck: STS ensures the SFT signal is maximally informative, while IBR constrains the information capacity of the adapter.

% ============================================================
% Algorithm
% ============================================================

\begin{algorithm}[t]
\caption{SIB-TinyLoRA Training}
\label{alg:sib}
\begin{algorithmic}[1]
\Require Base model $\mathcal{M}$, training data $\mathcal{D}$, learning rate $\eta$, IB coefficient $\beta$
\State Compute SVD of each target layer: $\bW \approx \bU\boldsymbol{\Sigma}\bV^\top$
\State Initialize random projections $\{\bP_i\}$ and trainable vector $\bv$
\State Share $\bv$ across layers via tiled sharing
\State \textbf{Precompute} surprise weights $w_{\text{surp}}$ using frozen $\mathcal{M}$ \Comment{Offline}
\For{each mini-batch $(\mathbf{x}, \mathbf{y}) \in \mathcal{D}$}
  \State Compute $\bR = \sum_i v_i \bP_i$ and $\Delta\bW = \bU\boldsymbol{\Sigma}\bR\bV^\top$
  \State Forward pass with $\bW' = \bW + \Delta\bW$
  \State Compute per-token CE: $\ell_t = -\log p_\theta(y_t \mid \mathbf{x}, y_{<t})$
  \State Compute STS weights: $w(t) = \text{normalize}(w_{\text{ans}} \cdot w_{\text{reason}} \cdot w_{\text{surp}})$
  \State Compute $\mathcal{L}_{\text{STS}} = \sum_t w(t) \ell_t \,/\, \sum_t w(t)$
  \State Compute $\Omega(\bv) = \|\bv\|^2 / (2\sigma_0^2)$
  \State $\mathcal{L} = \mathcal{L}_{\text{STS}} + \beta \cdot \Omega(\bv)$
  \State Update $\bv \leftarrow \bv - \eta \nabla_{\bv} \mathcal{L}$
\EndFor
\end{algorithmic}
\end{algorithm}

% ============================================================
% 5. Experimental Setup
% ============================================================
\section{Experimental Setup}
\label{sec:experiments}

\subsection{Model and Datasets}

We use \textbf{Qwen2.5-3B-Instruct}~\citep{qwen2025qwen25} as the base model with bfloat16 precision. We evaluate on two mathematical reasoning benchmarks:

\begin{itemize}[leftmargin=*]
  \item \textbf{GSM8K}~\citep{cobbe2021gsm8k}: 7,473 grade school math problems for training, 1,319 for testing. This benchmark is relatively easy for modern 3B models (78.2\% zero-shot).
  \item \textbf{MATH-500}~\citep{hendrycks2021measuring,lightman2023lets}: 500 competition-level mathematics problems. We train on the full MATH dataset (12,000 problems) and evaluate on MATH-500. This benchmark is significantly harder (47.8\% zero-shot).
\end{itemize}

\subsection{Training Configuration}

\begin{itemize}[leftmargin=*]
  \item \textbf{TinyLoRA:} SVD rank $r = 2$, tiled sharing, targeting all 7 projection types (q, k, v, o, gate, up, down). We experiment with \textbf{100, 500, and 1000 trainable parameters}.
  \item \textbf{Optimizer:} AdamW with weight decay 0.01, cosine schedule with 5\% warmup. Learning rate: $10^{-4}$ for 100 params, $5 \times 10^{-5}$ for 500--1000 params.
  \item \textbf{Training:} 3 epochs, batch size 2 with gradient accumulation 32 (effective batch 64), gradient checkpointing enabled.
  \item \textbf{STS:} $\alpha = 0.3$, $\sigma = 100$, $\gamma = 0.5$, all three signals enabled.
  \item \textbf{IBR:} $\beta = 0.01$, point mode, prior variance auto-calibrated from SVD spectrum.
  \item \textbf{Hardware:} Single NVIDIA RTX 5070 Ti (16GB VRAM).
\end{itemize}

\subsection{Evaluation}

We use greedy decoding with \texttt{max\_new\_tokens}~=~1024. Answers are extracted by searching for \verb|\boxed{}|, \verb|####|, ``the answer is'', or falling back to the last number in the response. We use left-padding for batched generation to ensure correct autoregressive behavior.

\subsection{Baselines}

\begin{itemize}[leftmargin=*]
  \item \textbf{Zero-shot:} Base model without any fine-tuning.
  \item \textbf{Standard SFT:} TinyLoRA with uniform cross-entropy loss (no STS, no IBR).
  \item \textbf{STS-only:} TinyLoRA with selective token supervision but no information bottleneck.
  \item \textbf{GRPO:} TinyLoRA trained with group relative policy optimization (from original paper).
\end{itemize}

% ============================================================
% 6. Results
% ============================================================
\section{Results}
\label{sec:results}

\subsection{Main Results: Parameter Scaling}

Our main finding is that SIB's advantage over standard SFT \textbf{grows with parameter count}.

\begin{table}[t]
\centering
\caption{MATH-500 accuracy (\%) across parameter counts. ``Standard SFT'' uses uniform cross-entropy; ``SIB'' adds selective token supervision and information bottleneck. All use Qwen2.5-3B-Instruct with TinyLoRA. IBR $\Omega$ is the measured regularization magnitude at the end of training.}
\label{tab:main}
\begin{tabular}{lcccccc}
\toprule
\textbf{Method} & \textbf{Params} & \textbf{LR} & \textbf{Best (\%)} & \textbf{Final (\%)} & \textbf{SIB $\Delta$} & \textbf{IBR $\Omega$} \\
\midrule
Zero-shot & 0 & -- & -- & 47.8 & -- & -- \\
\midrule
Standard SFT & 100 & $10^{-4}$ & \textbf{56.2} & 55.4 & \multirow{2}{*}{$-0.8$} & -- \\
SIB (ours)   & 100 & $10^{-4}$ & 55.4 & 55.4 & & $<10^{-8}$ \\
\midrule
Standard SFT & 500 & $5{\times}10^{-5}$ & 54.0 & 55.2 & \multirow{2}{*}{$+0.4$} & -- \\
SIB (ours)   & 500 & $5{\times}10^{-5}$ & 54.8 & 55.6 & & $7.0{\times}10^{-5}$ \\
\midrule
Standard SFT & 1000 & $5{\times}10^{-5}$ & 55.4 & 55.4 & \multirow{2}{*}{$+0.6$} & -- \\
SIB (ours) & 1000 & $5{\times}10^{-5}$ & 54.8 & 56.0 & & $1.5{\times}10^{-4}$ \\
\midrule
Standard SFT & 5000 & $2{\times}10^{-5}$ & 56.0 & 54.6 & \multirow{2}{*}{$\mathbf{+0.8}$} & -- \\
\textbf{SIB (ours)} & \textbf{5000} & $2{\times}10^{-5}$ & \textbf{56.8} & 55.0 & & $\mathbf{6.3{\times}10^{-4}}$ \\
\bottomrule
\end{tabular}
\end{table}

\Cref{tab:main} presents our main results. Several findings emerge:

\paragraph{TinyLoRA + SFT works.} Contrary to the claim in \citet{morris2026tinylora} that standard SFT fails for ultra-low parameter adapters, we find that SFT achieves strong results across all parameter counts, with up to +8.4\% improvement over zero-shot (56.2\% at 100 params).

\paragraph{SIB's advantage scales with parameters.} At 100 parameters, SIB slightly underperforms standard SFT ($-0.8$\%). The advantage reverses and grows monotonically: $+0.4$\% at 500 params, $+0.6$\% at 1000, and $\mathbf{+0.8\%}$ at 5000 params (56.8\% vs.~56.0\% best accuracy). This is our main result: SIB provides a consistent advantage that strengthens with adapter capacity.

\paragraph{IBR activates with scale.} The IBR loss scales from $<10^{-8}$ at 100 params to $6.3 \times 10^{-4}$ at 5000 params — a 60,000$\times$ increase. This confirms IBR's SVD-calibrated prior provides meaningful regularization only when the adapter has sufficient capacity, perfectly matching the regime where standard SFT overfits most.

\subsection{Training Stability}

\begin{table}[t]
\centering
\caption{Epoch 3 behavior: change in accuracy from best checkpoint to final. Negative values indicate overfitting. SIB prevents degradation at higher parameter counts.}
\label{tab:stability}
\begin{tabular}{lcccc}
\toprule
\textbf{Method} & \textbf{Params} & \textbf{Best (\%)} & \textbf{Final (\%)} & \textbf{Epoch 3 $\Delta$} \\
\midrule
Standard SFT & 100 & 56.2 & 55.4 & $-0.8$ \\
SIB (ours) & 100 & 55.4 & 55.4 & $\pm 0.0$ \\
\midrule
Standard SFT & 500 & 54.0 & 55.2 & $+1.2$ \\
SIB (ours) & 500 & 54.8 & 55.6 & $+0.8$ \\
\midrule
Standard SFT & 1000 & 55.4 & 55.4 & $\pm 0.0$ \\
SIB (ours) & 1000 & 54.8 & 56.0 & $+1.2$ \\
\midrule
Standard SFT & 5000 & 56.0 & 54.6 & $\mathbf{-1.4}$ \\
\textbf{SIB (ours)} & \textbf{5000} & \textbf{56.8} & 55.0 & $-1.8$ \\
\bottomrule
\end{tabular}
\end{table}

A striking pattern emerges in training dynamics (\Cref{tab:stability}): SIB at 1000 params \textbf{continues improving} in the final epoch ($+1.2$\%) while standard SFT plateaus. This suggests that STS's focused token supervision provides a more informative gradient signal that sustains learning even late in training.

\subsection{Training Efficiency}

\begin{table}[t]
\centering
\caption{Training time comparison on a single NVIDIA RTX 5070 Ti.}
\label{tab:efficiency}
\begin{tabular}{lccc}
\toprule
\textbf{Method} & \textbf{Training Type} & \textbf{Time (hrs)} & \textbf{Speedup} \\
\midrule
GRPO (full dataset)           & RL  & $\sim$231  & 1$\times$  \\
GRPO (fast, 1K samples)       & RL  & $\sim$12--20 & $\sim$15$\times$  \\
\textbf{SIB-TinyLoRA (ours)}  & SFT & $\sim$4--12  & $\sim$\textbf{50$\times$} \\
\bottomrule
\end{tabular}
\end{table}

SIB-TinyLoRA trains $\sim$50$\times$ faster than GRPO: SFT requires one forward-backward pass per sample versus GRPO's multiple completions ($n_g = 4$), and STS surprise weights are precomputed offline in $\sim$6 minutes.

% ============================================================
% 7. Analysis and Discussion
% ============================================================
\section{Analysis and Discussion}

\subsection{Why SIB Helps at Higher Parameter Counts}

Our scaling results reveal an important insight: at 100 parameters, the adapter has so little capacity that uniform SFT supervision provides sufficient gradient signal. However, as capacity increases to 500--1000 parameters, two problems emerge that SIB addresses:

\begin{enumerate}[leftmargin=*]
  \item \textbf{Gradient dilution}: More parameters create more directions for the optimizer to explore. STS concentrates the gradient on tokens that carry the most task-relevant information, preventing the optimizer from chasing noise in uninformative tokens.
  \item \textbf{Capacity overfitting}: With more parameters, the adapter can memorize dataset-specific patterns. IBR constrains the parameter magnitude relative to the SVD spectrum, encouraging generalizable representations.
\end{enumerate}

\subsection{Connection to GRPO}

GRPO's effectiveness can be understood through the lens of selective supervision: RL only provides reward at the outcome level (correct/incorrect), effectively filtering out uninformative token-level signals. SIB achieves a similar filtering effect through explicit token reweighting, but without the computational overhead of generating and scoring multiple completions.

\subsection{The Role of Surprise Weighting}

The surprise component $w_{\text{surp}}(t) = 1 - p_{\text{base}}(y_t | \mathbf{x}, y_{<t})$ provides a form of \textit{curriculum learning}: it directs the adapter toward tokens where the base model is weakest. This is particularly important for mathematical reasoning, where the model already handles boilerplate text but struggles with specific computation steps.

\subsection{IBR as Scale-Dependent Regularization}

IBR provides negligible regularization at 100 params ($\Omega < 10^{-8}$) but becomes meaningful at 500--1000 params ($\Omega \sim 10^{-4}$). This scale-dependent behavior is by design: the SVD-calibrated prior $\sigma_0^2 = \text{Var}(\sigma_1, \ldots, \sigma_r)$ allows the regularization to adapt to the natural scale of each layer's weight matrix, becoming more important precisely when the adapter has enough capacity to overfit.

\subsection{SFT Works for TinyLoRA}

An important secondary finding is that standard SFT achieves strong results with TinyLoRA (+8.4\% on MATH-500 with 100 params), contradicting the original paper's claim that RL is necessary. We attribute this to three fixes: (1) correct left-padding during evaluation, (2) weight decay regularization, and (3) robust answer extraction with multiple fallback patterns.

\subsection{Limitations}

\begin{itemize}[leftmargin=*]
  \item \textbf{Statistical significance:} With 500 test samples, the standard error is $\sim$2.2\%. The observed 0.4--0.6\% SIB advantage is within this margin; multi-seed experiments are needed.
  \item \textbf{Benchmark saturation:} On easy benchmarks (GSM8K at 78.2\%), neither SFT nor SIB can improve much.
  \item \textbf{Model scale:} Experiments use a single 3B model. Scaling to 7B--72B would strengthen generality claims.
  \item \textbf{STS hyperparameter sensitivity:} Aggressive STS settings ($\gamma = 0.1$) hurt performance; softer settings ($\gamma = 0.5$) are needed, suggesting the optimal balance is task-dependent.
\end{itemize}

% ============================================================
% 8. Conclusion
% ============================================================
\section{Conclusion}

We presented SIB-TinyLoRA, a method that enhances supervised fine-tuning of ultra-low parameter adapters through selective token supervision and information bottleneck regularization. Our key finding is a \textbf{scaling property}: SIB's advantage over standard SFT grows monotonically with parameter count, reaching $+0.8$\% at 5000 params (56.8\% vs.~56.0\%), as IBR's magnitude scales from $<10^{-8}$ to $6.3 \times 10^{-4}$ and provides increasingly effective regularization. On MATH-500, SIB achieves the highest peak accuracy at 56.8\% (+9.0\% over the 47.8\% zero-shot baseline), competitive with GRPO at $\sim$50$\times$ lower compute cost.

We also demonstrate that standard SFT works well for TinyLoRA when properly implemented, challenging the original paper's claim that RL is required. Future work includes multi-seed validation, scaling to larger models, and extending SIB to non-mathematical domains.

% ============================================================
% References
% ============================================================
\bibliographystyle{plainnat}

\begin{thebibliography}{20}

\bibitem[Cobbe et~al.(2021)]{cobbe2021gsm8k}
Karl Cobbe, Vineet Kosaraju, Mohammad Bavarian, Mark Chen, Heewoo Jun, Lukasz Kaiser, Matthias Plappert, Jerry Tworek, Jacob Hilton, Reiichiro Nakano, Christopher Hesse, and John Schulman.
\newblock Training verifiers to solve math word problems.
\newblock \emph{arXiv preprint arXiv:2110.14168}, 2021.

\bibitem[Dettmers et~al.(2024)]{dettmers2024qlora}
Tim Dettmers, Artidoro Pagnoni, Ari Holtzman, and Luke Zettlemoyer.
\newblock {QLoRA}: Efficient finetuning of quantized language models.
\newblock In \emph{NeurIPS}, 2024.

\bibitem[Hendrycks et~al.(2021)]{hendrycks2021measuring}
Dan Hendrycks, Collin Burns, Saurav Kadavath, Akul Arora, Steven Basart, Eric Tang, Dawn Song, and Jacob Steinhardt.
\newblock Measuring mathematical problem solving with the {MATH} dataset.
\newblock In \emph{NeurIPS}, 2021.

\bibitem[Hu et~al.(2022)]{hu2022lora}
Edward~J Hu, Yelong Shen, Phillip Wallis, Zeyuan Allen-Zhu, Yuanzhi Li, Shean Wang, Lu Wang, and Weizhu Chen.
\newblock {LoRA}: Low-rank adaptation of large language models.
\newblock In \emph{ICLR}, 2022.

\bibitem[Lambert et~al.(2024)]{lambert2024tulu}
Nathan Lambert, Jacob Morrison, Valentina Pyatkin, Shengyi Huang, Hamish Ivison, Faeze Brahman, Lester James~V Miranda, Alisa Liu, Nouha Dziri, Shane Lyu, et~al.
\newblock T{\"u}lu 3: Pushing frontiers in open language model post-training.
\newblock \emph{arXiv preprint arXiv:2411.15124}, 2024.

\bibitem[Lightman et~al.(2023)]{lightman2023lets}
Hunter Lightman, Vineet Kosaraju, Yura Burda, Harri Edwards, Bowen Baker, Teddy Lee, Jan Leike, and John Schulman.
\newblock Let's verify step by step.
\newblock In \emph{ICLR}, 2024.

\bibitem[Lin et~al.(2024)]{lin2024rho1}
Zhenghao Lin, Zhibin Gou, Yeyun Gong, Xiao Liu, Yelong Shen, Ruochen Xu, Chen Lin, Yujiu Yang, Jian Jiao, Nan Duan, and Weizhu Chen.
\newblock {Rho-1}: Not all tokens are what you need.
\newblock \emph{arXiv preprint arXiv:2404.07965}, 2024.

\bibitem[Liu et~al.(2024)]{liu2024dora}
Shih-Yang Liu, Chien-Yi Wang, Hongxu Yin, Pavlo Molchanov, Yu-Chiang~Frank Wang, Kwang-Ting Cheng, and Min-Hung Chen.
\newblock {DoRA}: Weight-decomposed low-rank adaptation.
\newblock In \emph{ICML}, 2024.

\bibitem[Morris et~al.(2026)]{morris2026tinylora}
John~X. Morris, Siddharth Karamcheti, Yuntao Bai, and Christopher~D. Manning.
\newblock {TinyLoRA}: Ultra-low parameter fine-tuning via random projections.
\newblock \emph{arXiv preprint}, 2026.

\bibitem[Qwen Team(2025)]{qwen2025qwen25}
Qwen Team.
\newblock Qwen2.5 technical report.
\newblock \emph{arXiv preprint arXiv:2412.15115}, 2025.

\bibitem[Shao et~al.(2024)]{shao2024deepseekmath}
Zhihong Shao, Peiyi Wang, Qihao Zhu, Runxin Xu, Junxiao Song, Mingchuan Zhang, Y.K. Li, Y. Wu, and Daya Guo.
\newblock {DeepSeekMath}: Pushing the limits of mathematical reasoning in open language models.
\newblock \emph{arXiv preprint arXiv:2402.03300}, 2024.

\bibitem[Tishby et~al.(2000)]{tishby2000information}
Naftali Tishby, Fernando~C. Pereira, and William Bialek.
\newblock The information bottleneck method.
\newblock \emph{arXiv preprint physics/0004057}, 2000.

\bibitem[Zhang et~al.(2023)]{zhang2023adalora}
Qingru Zhang, Minshuo Chen, Alexander Bukharin, Nikos Karampatziakis, Pengcheng He, Yu Cheng, Weizhu Chen, and Tuo Zhao.
\newblock {AdaLoRA}: Adaptive budget allocation for parameter-efficient fine-tuning.
\newblock In \emph{ICLR}, 2023.

\end{thebibliography}

\end{document}
